\section{Related Work}

\paragraph{Multi-Modal QA.} 
%
Multi-modal QA datasets have posed visual reasoning challenges for LMMs. 
Initially, the focus of benchmarks \cite{coco, MobasherParsVQACapsAB, yagcioglu-etal-2018-recipeqa, talmor2021multimodalqa, scienceqa, chang2022webqa, li2023seedbench, liu2023mmbench, yu2023mmvet} was on conducting QA tasks over simple images, primarily addressing questions such as understanding objects in an image and performing single-hop reasoning. 
%
Recently, more complex and nuanced benchmarks \cite{chen2022geoqa, lu2024mathvista} have emerged beyond the scope of understanding simple images to require complex mathematical reasoning over diagrams and plots. 
%
Beyond the scope of mathematical reasoning, MMMU \cite{yue2023mmmu} requires more complex visual reasoning in a diverse range of subjects including science, humanities, and engineering. 

\paragraph{Document QA.} 
%
Document QA is crucial in the field of NLP, focusing on extracting, synthesizing, and analyzing information from structured and unstructured documents. 
%
Early document QA benchmarks \cite{rajpurkar2016squad, bajaj2018msmarco, yang2018hotpotqa} involved short document QA, where questions were posed based on content from web pages such as those in Bing's search results or Wikipedia articles. 
%
Scientific paper QA benchmarks \cite{dasigi-etal-2021-dataset, pmlr-v202-lee23n} require LLMs to conduct multi-hop reasoning and long-context information processing. 
%
However, a notable gap exists in the integration of Multi-modal QA with Document QA, particularly in the context of scientific research, where it encompasses a blend of textual and visual data alongside complex textual information. 
%
\data, bridging this gap, is a benchmark for evaluating foundation models' abilities in both multi-modal and multi-document reasoning.

